\section{Exact implementations and accompanying data}
\label{app:implementation}

The finite proof uses integer arithmetic, finite sets, and exhaustive
search. The accompanying project contains two complete Python programs,
both using only the standard library. This appendix relates their
operations to the mathematical rules. The first program is reproduced
below in full; the second and the complete deduction traces are supplied
as separate files in the project.

\subsection{The compact implementation}

The function \texttt{forests} implements the recursion in
Lemma~\ref{finite:coverage}. Its nested function continues to extend a
tuple even when that tuple is not itself selected, so no valid longer
tuple is lost. An empty tuple consumes one clique vertex, which ensures
termination even at zero amount. The function \texttt{data} implements
Lemma~\ref{finite:labels}: its adjacency lists contain required edges,
and its endpoint lists give $T(x)$.

For a pair $x<y$, \texttt{forbidden} considers all subsets of old paths
not meeting $x$ or $y$ internally, and every permitted endpoint choice.
It tests endpoint degrees and increasing-edge reachability, which is
equivalent to connectivity by Lemma~\ref{finite:labels}. Its interval
endpoints are precisely those of~\eqref{finite:extension-interval}.
The functions \texttt{empty} and \texttt{initial} create separate sets
for the matrix entries, copy each initial constraint to its symmetric
position, and leave the diagonal empty.

The function \texttt{exact} adds every instance of the required-path
rule. For each starting vertex, its state is a pair $(U,x)$ consisting
of the visited vertex set, encoded by a bitmask, and the last vertex.
Initially $U$ contains just the start. Each transition adds a required
neighbour outside $U$. Induction on the number of transitions proves
that the states at step $\ell$ are exactly the visited-set and endpoint
pairs of simple paths of length $\ell$ from the start. Two paths with
the same state have the same allowed continuations, so merging them
loses no deduction. The function records $k-\ell$ at every resulting
endpoint, as required by Lemma~\ref{finite:required-path-rule}.
The starting vertex cannot be revisited, and reversing a path proves
symmetry. When the required graph has one extra vertex, the initial
copy of the old matrix fills only the old rows and columns; the new
entries begin empty.

The function \texttt{pack} computes
Lemma~\ref{finite:ball-bounds}, omits only the dominated radii described
in the main text, and uses~\eqref{finite:compatibility} as adjacency
between options. Its recursive function \texttt{find} is the
inclusion-and-exclusion search proved in Section~\ref{finite:packing-subsection}.
The order of the options affects the running time but not the result.
The special assignment \texttt{req[0]=\{0\}} handles two singleton
options; every other compatibility test uses a positive interval.

The function \texttt{bad} tests a required-edge conflict or a packing
contradiction. The function \texttt{attach} makes a fresh copy of all
required adjacency lists and adds either an edge or one vertex with
two incident edges. Finally, \texttt{check} performs the procedure of
Proposition~\ref{finite:procedure-soundness}. All trials in a round use
the old matrix; the list of new consequences is added only after the
trials finish. The outer loops are exactly~\eqref{finite:domain}.

\subsection{A separate implementation}

The file \texttt{verification/code/independent\_checker.py} imports no part of the
compact implementation. It generates the same patterns by first
partitioning $t$ into component orders, then taking positive
compositions of the amounts on each component, identifying reversal,
and choosing multisets of components of equal order. This enumerates
the same complete invariant proved in Lemma~\ref{finite:coverage}.

Old paths and their endpoint extensions are explicit vertex lists.
Candidate extensions are checked for intersection, and their endpoint
graph is tested by an undirected connectivity algorithm with degree
counts. For each positive outside parameter, the program tests the
optimality and cycle-closure inequalities directly. Forbidden parameters
are stored as integer flags on unordered endpoint pairs. Required
simple paths use sets of visited vertices, with the same state induction
as above, rather than integer bitmasks.

The packing search retains all radii and branches by base vertex: it
chooses one compatible radius at that base, or omits the base. Every
compatible family has one of these choices. An upper bound obtained
by summing the largest remaining weight at each base safely prunes
branches. The returned witness is checked for distinct bases, pairwise
compatibility, and total weight greater than $k$.
These different representations and search choices give a separate
implementation of the proved rules. Both complete runs give every row
of Table~\ref{finite:results}.

\subsection{Deduction traces and reproduction}

The file \texttt{verification/data/}\allowbreak\texttt{certificates.jsonl} contains one record for each
of the $3099$ parameter-pattern instances. Each record gives its
parameters, pattern, initial forbidden flags, every successful added
prohibition grouped by simultaneous round, and a final contradiction
witness. A packing witness records triples
\[
 (\text{base vertex},\ \text{radius},\ \text{independence bound}).
\]
Radius $-1$ denotes the singleton option, while radius $0$ denotes the
outside-neighbour set. A flag's bit in position $d$
records the prohibition of an outside path with $d$ internal vertices.
In particular, $d=0$ records a forbidden edge, not a singleton option.
All final witnesses in the supplied run are packing contradictions.

These records are deduction traces. Their initial constraints and
required-path consequences are recomputed by the full program; the
records are not offered as proof objects for an additional independent
minimal validator. The independent computational check is the separate
implementation just described, together with the proved meaning of
each inference. The trace makes its successful deductions available
for inspection and exact comparison.

From the project directory, the command
\begin{lstlisting}
python3 verification/code/verify.py --output /tmp/cycle-clique-check
\end{lstlisting}
runs both complete programs in separate processes and retains their
full output, errors, runtime information, implementation hashes and
generated data. The wrapper checks the complete range $5\le k\le17$,
all $42$ parameter pairs, exactly $3099$ distinct pattern instances,
zero unresolved cases, every printed row, and equality of the generated
trace with the supplied file. Only integer and finite-set operations
enter the mathematical decisions; wall-clock times are recorded solely
as execution metadata.

Use an ordinary Python~3 interpreter without optimization flags.
The independent program also accepts restricted ranges for diagnostic
runs, but their completion is not the full finite check. The wrapper
above always runs and checks the complete domain. Building the paper
and rerunning the calculations require no service, private repository,
or additional mathematical input; the project README gives the build
command and software requirements.

\subsection{The complete compact program}

\lstinputlisting[language=Python,basicstyle=\ttfamily\scriptsize]{../verification/code/original_checker.py}
